\section{Tipo de pesquisa e hipótese}
Informe qual tipo de pesquisa será conduzido (exploratória, empírica, formal etc.) e registre as hipóteses que serão avaliadas. Caso não haja hipótese formal, descreva o critério de sucesso esperado.

\section{Procedimentos da abordagem proposta}
A pesquisa bibliográfica foi conduzida de forma sistemática, estruturando-se em três temas principais: Hardware, Software e Redes. Para a seleção dos trabalhos, utilizaram-se as bases de dados Google Scholar, MDPI e Intechopen.

O processo de seleção seguiu um fluxo de filtragem (funil). Inicialmente, foram identificados 120 artigos por meio da análise de títulos, utilizando strings de busca majoritariamente em inglês, tais como ``iot AND energy efficiency'' e ``ESP32 AND camera monitoring''. Como ferramentas auxiliares na descoberta de conexões entre autores e na triagem inicial de relevância, foram utilizadas as plataformas ResearchRabbit e Gemini Deep Research.

Na etapa de refinamento, aplicou-se a leitura dos resumos (abstracts) e a verificação de disponibilidade integral dos textos. Foram adotados como critérios de inclusão: artigos publicados entre 2020 e 2026, relevância direta com sistemas de telemetria e foco em implementações práticas. Como critério de exclusão, descartaram-se trabalhos com acesso restrito ou que não apresentassem detalhes técnicos sobre a integração dos eixos propostos.

\begin{table}[H]
\centering
\caption{Strings de busca utilizadas no levantamento bibliográfico.}
\label{tab:strings_busca}
\begin{tabular}{@{}lll@{}}
\toprule
\textbf{ID} & \textbf{Eixo Temático} & \textbf{String de Busca} \\ \midrule
ST01 & Redes & "IoT" AND "\gls{lorawan}" AND "water waste" \\
ST02 & Contexto & smart water meter market \\
ST03 & Contexto & smart water meter costs and benefits \\
ST04 & Hardware & "IoT" AND "energy efficiency" \\
ST05 & Hardware & "ESP32" AND "camera monitoring" \\
ST06 & Hardware & "\gls{esp32cam}" AND "analog meter" AND "OCR" \\
ST08 & Software & "automatic meter reading" OCR accuracy "error rate" evaluation \\
ST09 & Redes & "packet delivery ratio" \gls{lorawan} "latency" monitoring system \\
ST010 & Contexto & "water consumption prediction" "mean absolute error" forecasting metrics \\ \bottomrule
\end{tabular}
\end{table}

\section{Ferramentas, tecnologias e ambientes}
Liste softwares, linguagens, bibliotecas, frameworks e ambientes (laboratórios, nuvem, simulação) que serão usados. Justifique a escolha com base no problema e nas métricas.

\section{Procedimentos experimentais e métricas}
Para avaliar a eficácia do sistema de telemetria proposto e garantir a objetividade dos resultados, serão adotados indicadores quantitativos divididos em três dimensões críticas: desempenho de rede, precisão de visão computacional, eficiência energética e acurácia preditiva.

\subsubsection{Desempenho da Rede de Comunicação}
A robustez da transmissão via \gls{lorawan} será quantificada através da Taxa de Entrega de Pacotes (\gls{pdr}). 
Conforme a metodologia de \citeonline{bravo2025ants}, o PDR é definido pela razão entre o número de pacotes recebidos pelo \textit{gateway} e o número de pacotes enviados pelo nó sensor:

\begin{equation}
\label{eq:pdr}
PDR = \left( \frac{\text{Pacotes Recebidos}}{\text{Pacotes Enviados}} \right) \times 100
\end{equation}

Além do \gls{pdr}, será medida a latência de ponta a ponta, compreendendo o tempo decorrido desde o disparo da captura da imagem no ESP32 até a disponibilidade do dado processado no \textit{dashboard}.

% =================================================================================================

\subsubsection{Acurácia da Visão Computacional (OCR)}
A validação do algoritmo de \gls{ocr}, responsável pela extração dos dígitos dos hidrômetros analógicos, será realizada sob diferentes condições de iluminação e ângulos de captura, visando simular cenários reais de instalação. 

Conforme discutido por \citeonline{barkat2024smart}, a eficácia de modelos de \textit{Deep Learning} aplicados ao monitoramento de água deve ser medida pela sua capacidade de generalização. 
Para este trabalho, a métrica principal será a Acurácia de Leitura ($Acc_{ocr}$), definida pela razão entre as predições que coincidem exatamente com o valor real (verdade de campo) e o total de amostras testadas:

\begin{equation}
Acc_{ocr} = \frac{VP + VN}{VP + VN + FP + FN}
\end{equation}

Onde $VP$, $VN$, $FP$ e $FN$ representam, respectivamente, os Verdadeiros Positivos, Verdadeiros Negativos, Falsos Positivos e Falsos Negativos obtidos através da Matriz de Confusão. 
Esta abordagem quantitativa permite identificar se erros de leitura estão associados a fatores externos, como reflexos na cúpula de vidro do medidor ou distorções de perspectiva.

% =================================================================================================

\subsubsection{Autonomia Energética}
A eficiência do hardware será validada através da medição do consumo de corrente em diferentes estados de operação (\textit{Deep Sleep}, captura de imagem e transmissão). 
A autonomia estimada ($A$) em dias será calculada pela relação entre a capacidade da bateria ($C_{bat}$) e o consumo médio diário ($I_{médio}$):

\begin{equation}
A = \frac{C_{bat}}{I_{médio} \times 24}
\end{equation}

% Explique como os experimentos serão configurados, quais variáveis serão controladas e quais métricas demonstram que a proposta cumpre os objetivos (ex.: acurácia, latência, custo). Inclua tabelas que antecipem o plano de coleta, por exemplo:
% \begin{table}[htb]
%   \centering
%   \caption{Plano de experimentação}
%   \begin{tabular}{lccc}
%     \toprule
%     Experimento & Variável controlada & Métrica principal \\
%     \midrule
%     Comparação com baseline & algoritmo & acurácia \\
%     Escalabilidade & carga & tempo médio \\
%     \bottomrule
%   \end{tabular}
%   \legend{Exemplo de como documentar experimentos.}
% \end{table}

% \textit{Propor algo não basta: demonstre como será validado.}

